% =============================================================================
% Section 12: MCQ Experiments
% Bib: refs/mcq_experiments.bib
%
% ACM-format paper fragment. The \subsection blocks below sit under the paper's
% MCQ "ATLAS Recreation" \section{} (see 10_mcq_atlas_recreation.tex). No
% \documentclass or \section here; \input this file after section 11.
%
% FIGURES REFERENCED (\includegraphics{figures/<basename>}).
% Integrator: each basename below was copied into ./figures/ from the ABSOLUTE
% source path shown.
%   figures/trainsize_corr_ext.png
%     <- /Users/arhant/Documents/EDLM/olmo-eval-full/AdaptiveTesting/Research/01_MCQ_ATLAS/figures/trainsize_corr_ext.png
%   figures/selection_curve_combined.png
%     <- /Users/arhant/Documents/EDLM/olmo-eval-full/AdaptiveTesting/Research/01_MCQ_ATLAS/data/model_diversity_selection/selection_curve_combined.png
%   figures/param_extrapolation_combined.png
%     <- /Users/arhant/Documents/EDLM/olmo-eval-full/AdaptiveTesting/Research/01_MCQ_ATLAS/data/param_extrapolation/param_extrapolation_combined.png
%   figures/calib_range_to_7b_r_A_vs_B.png
%     <- /Users/arhant/Documents/EDLM/olmo-eval-full/AdaptiveTesting/Research/01_MCQ_ATLAS/data/calib_range_to_7b/calib_range_to_7b_r_A_vs_B.png
%   figures/size_balance_r_comparison.png
%     <- /Users/arhant/Documents/EDLM/olmo-eval-full/AdaptiveTesting/Research/01_MCQ_ATLAS/data/size_balanced_recal/size_balance_r_comparison.png
%   figures/se_sweep_corr_vs_se.png
%     <- /Users/arhant/Documents/EDLM/olmo-eval-full/AdaptiveTesting/Research/01_MCQ_ATLAS/figures/se_sweep_corr_vs_se.png
%   figures/se_sweep_items_vs_se.png
%     <- /Users/arhant/Documents/EDLM/olmo-eval-full/AdaptiveTesting/Research/01_MCQ_ATLAS/figures/se_sweep_items_vs_se.png
%   figures/compare_2pl_vs_3pl_se0.3.png
%     <- /Users/arhant/Documents/EDLM/olmo-eval-full/AdaptiveTesting/Research/01_MCQ_ATLAS/figures/compare_2pl_vs_3pl_se0.3.png
%
% TABLE SOURCE DATA (absolute):
%   Table (train size)       <- .../01_MCQ_ATLAS/data/trainsize_corr.csv, trainsize_corr_ext.csv
%   Table (model selection)  <- .../01_MCQ_ATLAS/data/model_diversity_selection/README.md, results.csv
%   Table (range restriction)<- .../01_MCQ_ATLAS/data/calib_range_to_7b/README.md, results_summary.csv
%   Table (out-of-range)     <- .../01_MCQ_ATLAS/data/param_extrapolation/README.md, param_extrapolation_results.csv
%   Table (size vs range)    <- .../01_MCQ_ATLAS/data/size_balanced_recal/README.md, summary.csv
%   Table (stopping rule)    <- .../01_MCQ_ATLAS/data/se_sweep/se_sweep_combined.csv
%   Table (1PL/2PL/3PL)      <- .../se_sweep/se_sweep_combined.csv,
%                               .../data/compare_2pl_vs_3pl_se0.3.csv (GPQA; r derived),
%                               .../data/pedagogy_feasibility/rasch_1pl/pedagogy_1pl_vs_2pl_summary.csv
%
% PREAMBLE REQUIREMENTS (main document): \usepackage{graphicx}, \usepackage{xcolor},
%   \usepackage{booktabs}. If \url appears in the .bib, load \usepackage{hyperref} or {url}.
%
% Convention: the outline PDF gives only bullet headings for these experiments,
%   so the prose here is AI-added and wrapped in \aiadd{...} (renders blue).
%   Each fully-AI \subsection closes with an \aibullets{...} takeaway list.
%   [CITATION NEEDED] and pending-result markers render red.
% =============================================================================

\providecommand{\aiadd}[1]{\textcolor{blue}{#1}}
\providecommand{\aibullets}[1]{{\color{blue}#1}}

\aiadd{The experiments below measure what an ATLAS-style item bank costs to
build and what it returns. Every run uses the same validated calibration and
adaptive-testing pipeline~\cite{li2025atlas}: a chunked 3PL fit in the R
\texttt{mirt} package~\cite{chalmers2012mirt}, mean-sigma chunk linking, a
positive-discrimination ($a>0$) item filter, and an EAP-ability, Fisher-information
p-IRT computerized adaptive test~\cite{lord1980irt,wainer2000cat}. Unless stated
otherwise, the reported correlation is the Pearson $r$ between the CAT-predicted
benchmark accuracy and the true full-benchmark accuracy on held-out models, and
the stopping target is a posterior standard error of $0.3$. The five questions
are: how many calibration models are needed, whether a good choice of models
beats a random one, how the calibration parameter range constrains transfer,
where to set the stopping rule, and when the discrimination and guessing
parameters are worth their calibration cost.}

% =============================================================================
\subsection{Grounding the hypothesis: how many calibration models}
\label{subsec:mcq-trainsize}

\aiadd{IRT calibration estimates the item parameters from the response patterns
of the calibration models~\cite{lord1980irt,embretson2000irt}. A held-out model's
ability estimate then carries two sources of error: the posterior error from the
items it answered, and the uncertainty in the item parameters themselves. The
second term shrinks as the calibration pool grows, and it feeds directly into the
reported ability error~\cite{tsutakawa1990effect,patton2013influence}. The first practical
question is therefore how many calibration models a benchmark needs before its
bank is stable.}

\aiadd{We sweep the calibration-pool size on ARC-Challenge~\cite{clark2018arc},
refitting the 3PL bank on $N$ models and scoring the same 13-model held-out set
at each $N$ (Table~\ref{tab:mcq-trainsize}, Figure~\ref{fig:mcq-trainsize}).
Below about 15 models the correlation is unusable and swings between $0.45$ and
$0.92$ from one $N$ to the next, because a 3PL fit on so few respondents against
several hundred items is badly over-parameterized. From about 25 models up the
correlation settles into the $0.85$ to $0.94$ band: $r=0.938$ at $N=25$, $0.926$
at $N=45$, $0.898$ at $N=50$, and $0.846$ at $N=60$. The remaining
point-to-point movement is calibration noise, not a trend, so single points in
this range should be read as a band rather than a ranking.}

\aiadd{Two details matter for planning. First, the adaptive test stays about 4 to
5 items long at every pool size, so more calibration models buy calibration
stability, not shorter diagnostics. Second, a five-shuffle run across three MCQ
benchmarks reproduces the same shape, reaching a mean $r$ at $N=60$ of $0.897$
on ARC-Challenge, $0.845$ on SciQ, and $0.806$ on OpenBookQA. About 25 to 40
diverse models is enough to hold correlation near $0.85$ to $0.9$ on these
benchmarks, and below roughly 15 models the correlation collapses.}

\begin{table}[t]
  \centering
  \caption{\aiadd{Calibration-pool size against held-out correlation on
  ARC-Challenge (13 held-out models, $\mathrm{SE}\le0.3$). ``Items kept'' is the
  bank size after the $a>0$ filter; it grows with $N$ as more items show response
  variance. ``CAT items'' is the mean administered per model. The $N=60$ point
  ($r=0.846$) is reported in the text from the extended run. Source:
  \texttt{trainsize\_corr.csv}.}}
  \label{tab:mcq-trainsize}
  \small
  \begin{tabular}{r r r r r}
    \toprule
    $N$ models & Items kept & CAT items & Pearson $r$ & MAE \\
    \midrule
    5   & 444 & 4.5 & 0.555 & 0.104 \\
    10  & 589 & 4.1 & 0.454 & 0.096 \\
    15  & 628 & 4.2 & 0.919 & 0.047 \\
    20  & 657 & 4.1 & 0.764 & 0.075 \\
    25  & 668 & 4.7 & 0.938 & 0.044 \\
    30  & 678 & 4.9 & 0.856 & 0.057 \\
    35  & 716 & 5.5 & 0.905 & 0.043 \\
    40  & 733 & 5.5 & 0.830 & 0.052 \\
    45  & 752 & 5.4 & 0.926 & 0.034 \\
    50  & 759 & 5.2 & 0.898 & 0.042 \\
    \bottomrule
  \end{tabular}
\end{table}

\begin{figure}[t]
  \centering
  \includegraphics[width=\columnwidth]{figures/trainsize_corr_ext.png}
  \caption{\aiadd{Held-out correlation against calibration-pool size on
  ARC-Challenge at $\mathrm{SE}\le0.3$. The curve is noisy below about 15 models,
  then stabilizes in the $0.85$ to $0.94$ band from about 25 models up. Source:
  \texttt{figures/trainsize\_corr\_ext.png}.}}
  \label{fig:mcq-trainsize}
\end{figure}

\aibullets{%
\begin{itemize}
  \item Correlation is noisy below about 15 calibration models ($r$ swings $0.45$
  to $0.92$) and stabilizes at $0.85$ to $0.94$ from about 25 to 40 models up.
  \item The adaptive test stays about 4 to 5 items long at every pool size, so
  calibration models buy stability rather than shorter tests.
  \item Small-$N$ 3PL fits are over-parameterized, so single points below 15
  models are noise; read the band, not the individual value.
\end{itemize}}

% =============================================================================
\subsection{Model diversity against random selection}
\label{subsec:mcq-diversity}

\aiadd{If correlation stabilizes once the pool is diverse enough, then choosing
the calibration models well should reach the same correlation as a larger random
draw. We test this with the same pipeline, changing only which $N$ models enter
calibration, and scoring every strategy on one fixed seed-7 held-out split so the
curves are comparable~\cite{fourrier2024openllm}. Three strategies compete:
\emph{random} (five seeds, reported as mean and spread), \emph{ability spread}
(models chosen to cover the accuracy range evenly), and \emph{response diversity}
(a farthest-point selection on the model-by-item correctness matrix under Hamming
distance).}

\aiadd{Choosing models well reaches the random plateau with fewer models
(Table~\ref{tab:mcq-diversity}, Figure~\ref{fig:mcq-diversity}). On
IFEval~\cite{zhou2023ifeval} random saturates at $r=0.919$ near $N=120$; ability
spread matches that plateau at $N=80$, where random would need about 113 models,
a $1.41\times$ saving. On MATH~\cite{hendrycks2021math} random plateaus at
$r=0.865$ near $N=120$; response diversity matches the plateau with about
$1.59\times$ fewer models, and at $N=50$ it already reaches $r=0.889$, above the
120-model random plateau. In other words, 50 well-chosen models can beat a
random pool more than twice that size.}

\aiadd{The honest caveat is the same as in
Section~\ref{subsec:mcq-trainsize}: at the small-$N$ end these are 3PL fits on 10
to 30 respondents against 500 to 1{,}000 items, so the banks are over-parameterized
and shrink as constant items drop. The smallest-$N$ points are noisy and should be
read as a trend. The smart strategies are also deterministic, so they show a single
curve with no seed band, while random carries the spread across five seeds. The
savings factors come from interpolating the random mean curve to the match
threshold, so they are estimates rather than exact counts.}

\begin{table}[t]
  \centering
  \caption{\aiadd{Reaching the random-selection plateau with fewer models
  ($\mathrm{SE}\le0.3$). Each smart strategy is listed with the pool size at which
  it first matches the random plateau, the correlation there, and the
  interpolated model saving against random. Source:
  \texttt{model\_diversity\_selection/README.md} and \texttt{results.csv}.}}
  \label{tab:mcq-diversity}
  \small
  \begin{tabular}{l l r r r}
    \toprule
    Benchmark & Strategy & $N$ & Pearson $r$ & Saving \\
    \midrule
    IFEval & random plateau    & 120 & 0.919 & n/a       \\
    IFEval & ability spread    & 80  & 0.921 & 1.41$\times$ \\
    IFEval & response diversity& 120 & 0.912 & 0.94$\times$ \\
    \midrule
    MATH   & random plateau    & 120 & 0.865 & n/a       \\
    MATH   & ability spread    & 50  & 0.859 & 0.95$\times$ \\
    MATH   & response diversity& 30  & 0.837 & 1.59$\times$ \\
    \bottomrule
  \end{tabular}
\end{table}

\begin{figure}[t]
  \centering
  \includegraphics[width=\columnwidth]{figures/selection_curve_combined.png}
  \caption{\aiadd{Held-out correlation against pool size for random, ability-spread,
  and response-diversity selection on IFEval and MATH ($\mathrm{SE}\le0.3$). The
  shaded band is the random seed spread; markers show where a smart strategy first
  reaches the random plateau. Source: \texttt{figures/selection\_curve\_combined.png}.}}
  \label{fig:mcq-diversity}
\end{figure}

\aibullets{%
\begin{itemize}
  \item A diversity-aware or ability-spread choice reaches the random plateau with
  $1.4\times$ to $1.6\times$ fewer models (IFEval ability spread $1.41\times$; MATH
  response diversity $1.59\times$).
  \item On MATH, 50 response-diverse models reach $r=0.889$, above the $r=0.865$
  plateau of a 120-model random pool.
  \item Small-$N$ 3PL fits are noisy and over-parameterized, so the gains should be
  read from the trend, not from single low-$N$ points.
\end{itemize}}

% =============================================================================
\subsection{Oracle latent-ability subset selection}
\label{subsec:mcq-oracle}

\textcolor{red}{[ORACLE SUBSET RESULT PENDING: theoretical best-subset selection
reverse-engineered from final calibration into latent-ability groups; $\sim$300 initial
models, 50-model steps; fill on completion.]}

% =============================================================================
\subsection{Parameter range: restricting the pool and testing out of range}
\label{subsec:mcq-range}

\aiadd{A team calibrating for a target model size wants to know whether it can use
a narrow, matched calibration pool, and whether a bank calibrated on one size
range can score models in another. We run three experiments that pull in different
directions at first and then agree on one rule: what drives correlation is the
ability spread of the calibration pool, not how close its parameter range sits to
the models under test.}

\aiadd{\textbf{Reducing models by specifying a range.} We fix the test set to 7B
models (parameters $\ge 6.5$B) and compare two matched-size calibration pools:
pool A drawn from $[0,5)$B, which sits entirely below the test range, and pool B
drawn from $[3,6.5)$B, which sits right beneath it. Proximity would favor B. It
does not (Table~\ref{tab:mcq-calibrange}, Figure~\ref{fig:mcq-calibrange}). The
wider low pool A wins on every benchmark at $\mathrm{SE}\le0.3$: $r=0.961$ against
$0.831$ on MATH, $0.907$ against $0.871$ on IFEval, $0.697$ against $0.512$ on
BBH~\cite{suzgun2023bbh}, and a near-tie on GPQA~\cite{rein2024gpqa}
($0.632$ against $0.594$). The pool that spans a wider ability range predicts 7B
models better than the pool that sits closest to them.}

\aiadd{\textbf{Using out-of-range models.} We then remove the middle of the range
entirely: calibrate on a LOW band ($\le 2$B) or a HIGH band ($\ge 3.5$B) with the
$(2,3.5)$B middle excluded, and test across the gap, using in-range held-out
controls for comparison (Table~\ref{tab:mcq-extrap},
Figure~\ref{fig:mcq-extrap}). A bank does extrapolate across the gap when the
benchmark links well. On IFEval the transfer is essentially free, with
out-of-range correlations of $0.917$ (up) and $0.919$ (down) against a $0.910$
in-range control. On MATH it works with a real cost, $0.789$ up and $0.859$ down.
The cost is asymmetric in every benchmark: calibrating on small models and testing
large (the up direction) is harder than the reverse, because a bank built on small
models is dominated by items that saturate for large models and lose
discrimination at the top. GPQA shows the same mechanism from the other side: it
links weakly and is near guessing for small models, so out-of-range transfer
breaks down to $r=0.389$ up and $0.540$ down. That is a linking failure, not a
property of parameter range.}

\aiadd{\textbf{Range restriction against size skew.} The first ATLAS
recreation showed that refitting the published bank on only the $0.5$ to $7$B
slice drops held-out ARC correlation from $0.830$ to $0.589$ at
$\mathrm{SE}\le0.2$. That slice is about $95.6\%$ 7B models, so the drop could be
range restriction or a size-skewed pool. We hold $N$ fixed and vary only the size
mix (Table~\ref{tab:mcq-sizebal}, Figure~\ref{fig:mcq-sizebal}). De-skewing the
pool from about $95\%$ 7B down to $33$ to $50\%$ 7B does not move the correlation
back toward $0.83$. The balanced pools land at $r\approx0.56$ to $0.59$, right on
top of the full-slice result and slightly below the skewed pools at the same $N$.
Pool count does not rescue it either: skewed $r\approx0.65$ holds at $N=72$ and
$N=174$, and the full slice gives $0.589$ at $N=1{,}686$. Balancing does lower the
accuracy MAE (about $0.11$ against $0.13$, since small models sharpen absolute
calibration), but it costs rank-correlation stability, with the two-band spread
rising to about $0.20$ as scarce sub-2B models add noise.}

\aiadd{Read together, the three results give one rule. Pool A beats pool B
because A covers a wider ability range; the extrapolation bank transfers because a
well-linked bank carries its discrimination across a parameter gap; and
rebalancing sizes inside the restricted $0.5$ to $7$B slice cannot help because the
missing ingredient is the wide ability spread of the full ATLAS pool, not the size
histogram. Calibration-pool ability spread is the main driver, and matching or
overlapping the target parameter range is secondary.}

\begin{table}[t]
  \centering
  \caption{\aiadd{Reducing models by parameter range, tested on 7B models
  ($\mathrm{SE}\le0.3$, matched $N$). Pool A is $[0,5)$B (below the test range);
  pool B is $[3,6.5)$B (just below it). Pool B uses its whole pool each seed, so
  its spread is near zero. Source: \texttt{calib\_range\_to\_7b/results\_summary.csv}.}}
  \label{tab:mcq-calibrange}
  \small
  \begin{tabular}{l r r r}
    \toprule
    Benchmark & A: $[0,5)$B $r$ & B: $[3,6.5)$B $r$ & $B-A$ \\
    \midrule
    BBH    & 0.697 & 0.512 & $-0.185$ \\
    GPQA   & 0.632 & 0.594 & $-0.038$ \\
    IFEval & 0.907 & 0.871 & $-0.036$ \\
    MATH   & 0.961 & 0.831 & $-0.130$ \\
    \bottomrule
  \end{tabular}
\end{table}

\begin{figure}[t]
  \centering
  \includegraphics[width=\columnwidth]{figures/calib_range_to_7b_r_A_vs_B.png}
  \caption{\aiadd{Recovery of 7B-model accuracy from a wide out-of-range pool A
  ($[0,5)$B) against a narrow near-range pool B ($[3,6.5)$B) at matched $N$. The
  wider pool wins on every benchmark. Source:
  \texttt{figures/calib\_range\_to\_7b\_r\_A\_vs\_B.png}.}}
  \label{fig:mcq-calibrange}
\end{figure}

\begin{table*}[t]
  \centering
  \caption{\aiadd{Out-of-range transfer across an excluded middle band
  (LOW~$\le2$B, HIGH~$\ge3.5$B; $\mathrm{SE}\le0.3$, mean over 3 seeds). LOW$\to$LOW
  and HIGH$\to$HIGH are in-range controls; LOW$\to$HIGH extrapolates up and
  HIGH$\to$LOW extrapolates down. Penalty is control minus extrapolation on the
  shared test set. Source: \texttt{param\_extrapolation/README.md} and
  \texttt{param\_extrapolation\_results.csv}.}}
  \label{tab:mcq-extrap}
  \small
  \begin{tabular}{l c c c c c c}
    \toprule
    Benchmark & LOW$\to$LOW & HIGH$\to$HIGH & LOW$\to$HIGH (up) & HIGH$\to$LOW (down) & Penalty up & Penalty down \\
    \midrule
    IFEval & 0.910 & 0.955 & 0.917 & 0.919 & $+0.038$ & $-0.009$ \\
    MATH   & 0.891 & 0.907 & 0.789 & 0.859 & $+0.118$ & $+0.032$ \\
    GPQA   & 0.854 & 0.835 & 0.389 & 0.540 & $+0.445$ & $+0.314$ \\
    \bottomrule
  \end{tabular}
\end{table*}

\begin{figure*}[t]
  \centering
  \includegraphics[width=0.9\textwidth]{figures/param_extrapolation_combined.png}
  \caption{\aiadd{Out-of-range transfer for IFEval, MATH, and GPQA. Each panel
  shows the two in-range controls and the two cross-gap conditions. IFEval
  transfers freely, MATH transfers with an asymmetric cost, and GPQA breaks down
  because it links weakly. Source: \texttt{figures/param\_extrapolation\_combined.png}.}}
  \label{fig:mcq-extrap}
\end{figure*}

\begin{table}[t]
  \centering
  \caption{\aiadd{Range restriction against size skew on ARC ($\mathrm{SE}\le0.2$,
  60 held-out models; balanced and skewed rows are mean over 5 seeds). De-skewing
  the $0.5$ to $7$B pool does not recover the wide-pool correlation. The first two
  rows are references from the ATLAS recreation. Source:
  \texttt{size\_balanced\_recal/README.md} and \texttt{summary.csv}.}}
  \label{tab:mcq-sizebal}
  \small
  \begin{tabular}{l r r r r}
    \toprule
    Pool & $N$ & \% 7B & $r$ & MAE \\
    \midrule
    Published ATLAS bank & \multicolumn{1}{c}{$\sim$1000s} & wide  & 0.830 & 0.147 \\
    Full $0.5$ to $7$B   & 1{,}686 & 95.6 & 0.589 & 0.147 \\
    Skewed               & 72   & $\sim$95 & 0.652 & 0.130 \\
    Balanced (3-band)    & 72   & 33   & 0.594 & 0.106 \\
    Skewed               & 174  & $\sim$95 & 0.651 & 0.139 \\
    Balanced (2-band)    & 174  & 50   & 0.560 & 0.114 \\
    \bottomrule
  \end{tabular}
\end{table}

\begin{figure}[t]
  \centering
  \includegraphics[width=\columnwidth]{figures/size_balance_r_comparison.png}
  \caption{\aiadd{Held-out ARC correlation for skewed against size-balanced pools
  at matched $N$, with the published wide-pool bank and the full $0.5$ to $7$B slice
  for reference. Balancing the size mix does not recover the wide-pool
  correlation. Source: \texttt{figures/size\_balance\_r\_comparison.png}.}}
  \label{fig:mcq-sizebal}
\end{figure}

\aibullets{%
\begin{itemize}
  \item A wide out-of-range pool ($[0,5)$B) predicts 7B models better than a narrow
  near-range pool ($[3,6.5)$B) on every benchmark, so ability spread beats
  proximity to the test range.
  \item A well-linked bank extrapolates across an excluded parameter band (IFEval
  almost free, MATH with a modest cost), but transfer is asymmetric: calibrating
  small and testing large costs more than the reverse.
  \item Weakly-linking GPQA fails to transfer out of range ($r=0.39$ up), a linking
  failure rather than a range effect.
  \item Rebalancing the size mix inside the restricted $0.5$ to $7$B slice does not
  recover the wide-pool correlation ($r\approx0.56$ to $0.59$ against $0.83$); the
  binding constraint is range restriction, not size skew.
\end{itemize}}

% =============================================================================
\subsection{Stopping rule: the standard-error target}
\label{subsec:mcq-sestop}

\aiadd{The adaptive test stops when a model's ability estimate reaches a target
posterior standard error, so that target trades test length against precision and
correlation~\cite{wainer2000cat}. We sweep the target from $0.5$ down to $0.12$
for three banks: the published ATLAS 3PL bank replayed on ATLAS's own ARC
responses, a locally-fit 2PL bank, and a locally-fit 1PL bank on the same ARC
pool (Table~\ref{tab:mcq-sestop}, Figures~\ref{fig:mcq-sestop-corr}
and~\ref{fig:mcq-sestop-items}). The three banks respond very differently, so the
right target depends on which bank is in use.}

\aiadd{The well-calibrated ATLAS 3PL bank is already near its ceiling at a loose
target. It reaches $r=0.905$ in about 8 items at $\mathrm{SE}\le0.5$, and
tightening all the way to $0.12$ raises correlation only to $0.931$ while the test
grows to about 47 items. Past $\mathrm{SE}\le0.3$ the extra items mostly buy
precision, not ranking accuracy. The local 1PL bank behaves the opposite way:
correlation climbs steeply with the target, from $0.729$ at $\mathrm{SE}\le0.5$
(about 14 items) to $0.976$ at $\mathrm{SE}\le0.3$ (about 43 items) and $0.987$ at
$\mathrm{SE}\le0.15$ (182 items), so precision and correlation improve together but
the test gets long. The local 2PL bank sits at the 8-item floor with $r\approx0.86$
until the target passes $0.2$, dips slightly there, then jumps to $0.965$ at
$\mathrm{SE}\le0.12$ once it is allowed about 71 items.}

\aiadd{The practical reading is to set the target from the accuracy requirement
rather than by default. A well-calibrated 3PL bank delivers about $0.90$
correlation at $\mathrm{SE}\le0.3$ in roughly 9 items, and tightening further is
mostly wasted on ranking. A 1PL or 2PL bank needs a tighter target to reach the
same correlation, and pays for it in items.}

\begin{table}[t]
  \centering
  \caption{\aiadd{Stopping target against mean items and correlation for three
  banks on ARC: the published ATLAS 3PL bank on ATLAS's own responses (650 items),
  a local 2PL bank, and a local 1PL bank (759 items each). Source:
  \texttt{se\_sweep/se\_sweep\_combined.csv}.}}
  \label{tab:mcq-sestop}
  \small
  \begin{tabular}{r r r r r r r}
    \toprule
    & \multicolumn{2}{c}{ATLAS 3PL} & \multicolumn{2}{c}{Local 2PL} & \multicolumn{2}{c}{Local 1PL} \\
    \cmidrule(lr){2-3}\cmidrule(lr){4-5}\cmidrule(lr){6-7}
    SE~$\le$ & items & $r$ & items & $r$ & items & $r$ \\
    \midrule
    0.50 & 8.2  & 0.905 & 8.0  & 0.862 & 13.7  & 0.729 \\
    0.30 & 9.2  & 0.906 & 8.0  & 0.862 & 43.4  & 0.976 \\
    0.20 & 13.0 & 0.915 & 14.6 & 0.842 & 100.5 & 0.977 \\
    0.15 & 32.2 & 0.925 & 32.9 & 0.920 & 182.0 & 0.987 \\
    0.12 & 47.2 & 0.931 & 71.1 & 0.965 & 303.5 & 0.985 \\
    \bottomrule
  \end{tabular}
\end{table}

\begin{figure}[t]
  \centering
  \includegraphics[width=\columnwidth]{figures/se_sweep_corr_vs_se.png}
  \caption{\aiadd{Correlation against the standard-error target for the ATLAS 3PL,
  local 2PL, and local 1PL banks on ARC. The 3PL bank is near ceiling at a loose
  target; the 1PL bank climbs steeply as the target tightens. Source:
  \texttt{figures/se\_sweep\_corr\_vs\_se.png}.}}
  \label{fig:mcq-sestop-corr}
\end{figure}

\begin{figure}[t]
  \centering
  \includegraphics[width=\columnwidth]{figures/se_sweep_items_vs_se.png}
  \caption{\aiadd{Mean items administered against the standard-error target for the
  same three banks. A tighter target costs far more items for the 1PL bank than for
  the 3PL bank. Source: \texttt{figures/se\_sweep\_items\_vs\_se.png}.}}
  \label{fig:mcq-sestop-items}
\end{figure}

\aibullets{%
\begin{itemize}
  \item The published 3PL bank reaches $r\approx0.90$ at $\mathrm{SE}\le0.3$ in
  about 9 items; tightening to $0.12$ adds items (about 47) for a small correlation
  gain (to $0.931$).
  \item The local 1PL bank trades length for correlation, rising from $r=0.729$ at
  $\mathrm{SE}\le0.5$ to $0.987$ at $\mathrm{SE}\le0.15$ as the test grows from 14
  to 182 items.
  \item Set the target from the accuracy requirement: a well-calibrated 3PL bank
  gains little past $\mathrm{SE}\le0.3$, while 1PL and 2PL banks need a tighter
  target and more items to match it.
\end{itemize}}

% =============================================================================
\subsection{Value of the discrimination and guessing parameters}
\label{subsec:mcq-1pl2pl3pl}

\aiadd{The 1PL/Rasch model fits only item difficulty~\cite{rasch1960}, the 2PL
adds a discrimination slope, and the 3PL adds a lower-asymptote guessing term
\cite{birnbaum1968}. Each extra parameter can sharpen the test, but only if the
calibration pool is large enough to estimate it reliably. The stopping-rule sweep
and two paired comparisons show when the extra parameters earn their cost
(Table~\ref{tab:mcq-param}, Figure~\ref{fig:mcq-2pl3pl}).}

\aiadd{The guessing parameter is decisive on a near-guessing benchmark. On GPQA,
which is four-choice and close to chance for small models, the 3PL bank recovers
held-out accuracy at $r=0.737$ in about 13 items, while the 2PL bank on the same
110 models collapses to $r=0.095$. Without a lower asymptote the 2PL cannot fit a
benchmark whose weak models cluster near the guessing floor.}

\aiadd{When calibration models are scarce, the discrimination slope becomes a
liability and the simpler 1PL is more robust. On the Pedagogy
benchmark~\cite{lelievre2025pedagogy}, calibrated on 40 models, the 1PL bank ranks
held-out models at $r=0.890$ at $\mathrm{SE}\le0.3$ against the 2PL bank's $0.716$.
The same ordering holds on the in-house ARC pool, where the 1PL bank reaches
$r=0.976$ at $\mathrm{SE}\le0.3$ against the 2PL bank's $0.862$. With few models
the estimated slopes are noisy and drag the ranking down.}

\aiadd{The discrimination slope pays off in efficiency once it is well estimated.
On ARC the 2PL bank hits $\mathrm{SE}\le0.3$ in 8 items against the 1PL bank's 43,
and it reaches $r=0.965$ at $\mathrm{SE}\le0.12$ in 71 items against the 1PL bank's
304. With a large pool the full 3PL gives both efficiency and accuracy: the
published ATLAS bank, calibrated on thousands of models, recovers ARC accuracy at
$r=0.906$ in about 9 items at $\mathrm{SE}\le0.3$. The reading is that 2PL and 3PL
add the most when calibration models are plentiful, and 1PL is the safer choice
when they are scarce, except on near-guessing benchmarks where the 3PL guessing
term is required regardless.}

\begin{table}[t]
  \centering
  \caption{\aiadd{Discrimination and guessing against calibration-pool size, all at
  $\mathrm{SE}\le0.3$. The near-guessing GPQA comparison shows the 3PL guessing term
  is required; the scarce-pool comparisons show 1PL ranks better than 2PL. Source:
  \texttt{se\_sweep\_combined.csv}, \texttt{compare\_2pl\_vs\_3pl\_se0.3.csv} (GPQA;
  Pearson $r$ derived from the per-model file), and
  \texttt{pedagogy\_1pl\_vs\_2pl\_summary.csv}.}}
  \label{tab:mcq-param}
  \small
  \begin{tabular}{l l r r}
    \toprule
    Bank & Pool (benchmark, models) & $r$ & Items \\
    \midrule
    3PL & ARC, thousands   & 0.906 & 9.2  \\
    3PL & GPQA, 110         & 0.737 & 13.0 \\
    2PL & GPQA, 110         & 0.095 & 18.1 \\
    2PL & ARC, in-house     & 0.862 & 8.0  \\
    1PL & ARC, in-house     & 0.976 & 43.4 \\
    2PL & Pedagogy, 40      & 0.716 & 10.2 \\
    1PL & Pedagogy, 40      & 0.890 & 43.8 \\
    \bottomrule
  \end{tabular}
\end{table}

\begin{figure}[t]
  \centering
  \includegraphics[width=\columnwidth]{figures/compare_2pl_vs_3pl_se0.3.png}
  \caption{\aiadd{Predicted against actual accuracy for the 3PL and 2PL banks on
  GPQA at $\mathrm{SE}\le0.3$. The 3PL guessing term recovers the ranking
  ($r=0.737$); the 2PL bank, with no lower asymptote, collapses near the guessing
  floor ($r=0.095$). Source: \texttt{figures/compare\_2pl\_vs\_3pl\_se0.3.png}.}}
  \label{fig:mcq-2pl3pl}
\end{figure}

\aibullets{%
\begin{itemize}
  \item The 3PL guessing term is required on near-guessing benchmarks: on GPQA the
  3PL bank reaches $r=0.737$ while the 2PL bank collapses to $r=0.095$.
  \item When calibration models are scarce, 1PL ranks better than 2PL (Pedagogy
  $0.890$ against $0.716$; in-house ARC $0.976$ against $0.862$ at
  $\mathrm{SE}\le0.3$).
  \item The discrimination slope buys item-efficiency once well estimated (ARC 2PL
  reaches its target in 8 items against the 1PL bank's 43), and a large-pool 3PL is
  both efficient and accurate ($r=0.906$ in about 9 items).
\end{itemize}}

% =============================================================================
\subsection{Results}
\label{subsec:mcq-results}

\aiadd{The five experiments give a consistent recipe for standing up an
ATLAS-style MCQ diagnostic. Get the calibration pool diverse enough, spend the
model budget on ability spread rather than range matching, and pick the model
family and stopping target from the calibration budget and the accuracy
requirement.}

\aibullets{%
\begin{itemize}
  \item About 25 to 40 calibration models stabilize correlation at $0.85$ to
  $0.94$; below about 15 the fit collapses, and extra models buy stability, not
  shorter tests.
  \item Choosing models by ability spread or response diversity reaches the random
  plateau with $1.4\times$ to $1.6\times$ fewer models.
  \item Calibration-pool ability spread drives correlation more than matching the
  test parameter range: a wide out-of-range pool beats a narrow near-range one, and
  rebalancing sizes inside a restricted range does not recover the wide-pool
  correlation.
  \item Set the standard-error target from the accuracy requirement; a
  well-calibrated 3PL bank gains little past $\mathrm{SE}\le0.3$.
  \item Use 1PL when calibration models are scarce, 2PL or 3PL when they are
  plentiful, and the 3PL guessing term whenever the benchmark is near guessing.
\end{itemize}}
